% chapters/08/latihan.tex
\bagian{Latihan dan Refleksi}

\begin{enumerate}
  \item Sebutkan tiga langkah pada susunan keputusan sebuah agen, lalu
        jelaskan apa yang dilakukan pada masing-masing langkah.

  \item Sebuah NPC berada pada posisi $(100, 100)$ dan pemain pada posisi
        $(400, 500)$. Hitunglah vektor menuju pemain, jaraknya, dan vektor
        kecepatannya apabila kecepatan NPC ditetapkan 120 piksel per detik.

  \item Mengapa urutan pengurangan pada perhitungan vektor menuju pemain
        penting? Apa yang terjadi apabila urutannya dibalik?

  \item Jelaskan mengapa memakai vektor jarak secara langsung sebagai kecepatan
        merupakan cara yang salah. Sebutkan tiga akibatnya bagi permainan.

  \item Sebuah NPC bergerak dengan kecepatan 240 piksel per detik pada laju 60
        bingkai per detik, dan jarak hentinya 20 piksel. Tentukan apakah NPC
        akan bergetar di sekitar pemain, beserta perhitungannya.

  \item Jelaskan mengapa keputusan ``berdiam'' dan ``berhenti karena sudah
        cukup dekat'' perlu dibedakan, walaupun kecepatannya sama-sama nol.

  \item Jelaskan apa yang dimaksud dengan \istilah{hysteresis} pada batas
        persepsi, dan gejala apa yang muncul apabila cara itu tidak dipakai.

  \item Perhatikan Gambar~\ref{fig:reaktif-perencanaan}. Jelaskan mengapa NPC
        pada panel pertama tidak melakukan kesalahan apa pun, walaupun
        perilakunya tampak bodoh.
\end{enumerate}

\begin{refleksibox}
\begin{enumerate}
  \item Pada langkah keenam praktikum, berapa lama NPC menekan bongkahan sebelum
        Anda menghentikan program? Menurut Anda, bagaimana pemain merasakan
        perilaku itu?
  \item Bandingkan hasil pengukuran Anda pada langkah ketiga dan keempat.
        Seberapa besar perbedaannya, dan apakah perbedaan itu terlihat mata
        tanpa pengukuran?
  \item Pada langkah kelima, jari-jari persepsi mana yang menghasilkan
        permainan yang paling menantang? Nilai itu lebih besar atau lebih kecil
        daripada dugaan awal Anda?
\end{enumerate}
\end{refleksibox}
